\section{Introduction}\label{sec:intro}

Chain-of-thought (CoT) prompting has emerged as a central technique for
eliciting multi-step reasoning in large language models~\cite{Wei2022}. By
instructing a model to produce intermediate reasoning steps before arriving at
a final answer, CoT prompting substantially improves performance on arithmetic,
commonsense, and symbolic reasoning tasks. Despite its empirical success,
the mathematical structure underlying CoT prompt composition remains poorly
understood. When multiple CoT prompts are applied in sequence---refining,
extending, or correcting previous reasoning---fundamental questions arise: When
does iterated prompting converge to a stable reasoning pattern? What algebraic
structure governs the composition of prompts? How can one predict whether
repeated application of a prompt will improve, degrade, or cycle through
reasoning states?

Recent work has begun to address these questions from several directions.
Geshkovski et~al.~\cite{Geshkovski2023} model transformer layers as a
one-parameter semigroup of flow maps on the space of probability measures over
the unit sphere, proving that tokens cluster to a single point as depth tends
to infinity. Kim et~al.~\cite{Kim2025} establish a universal approximation
theorem for prompt-induced hypothesis classes, showing that a fixed transformer
backbone can approximate any continuous function by varying only the prompt.
Yang et~al.~\cite{Yang2025} model LLMs as improvement operators and observe
empirically that iterative refinement saturates in accuracy. Badea
et~al.~\cite{Badea2025} prove that iterated holomorphic operator transforms
converge to idempotent projections under a peripheral fixed-point property. On
the semigroup-theoretic side, Gluck and Haase~\cite{Gluck2018} develop the
theory of the semigroup at infinity and the Jacobs--de~Leeuw--Glicksberg
(JdLG) decomposition for bounded operator semigroups, while Gerlach and
Gluck~\cite{Gerlach2017} establish convergence criteria for positive operator
semigroups via AM-compactness.

Despite this progress, no prior work constructs a semigroup-theoretic framework
for CoT prompt composition or establishes formal algebraic criteria for
convergence of iterated prompting. The present paper fills this gap.

\subsection*{Contributions}
We provide the first rigorous algebraic framework for CoT prompt composition as
a semigroup of operators on a metric space of reasoning trajectories. Our main
contributions are:

\begin{itemize}[leftmargin=*,itemsep=2pt,topsep=2pt]
\item We define \emph{linear CoT operators} $T(A,B)$ on a trajectory space
  $\mathcal{M}_n$ and prove that they form a monoid under composition, with an
  explicit closed-form composition formula
  (Theorem~\ref{thm:semigroup}).

\item We establish a sharp contractivity criterion: iterates $T^n$ converge in
  the strong operator topology if and only if the spectral radius
  $\rho(B)<1$, with convergence rate $O(\rho(B)^n)$
  (Theorem~\ref{thm:contractivity}).

\item We characterize idempotent CoT operators as precisely those satisfying
  $B^2=B$ and $BA=0$, and describe the fixed-point set explicitly
  (Theorem~\ref{thm:idempotent}).

\item We give a complete four-case classification of convergence
  behavior---contractive, projective, oscillatory, and
  divergent---determined by the spectrum of the current-state matrix~$B$
  (Theorem~\ref{thm:classification}).

\item We apply the JdLG decomposition to separate reasoning trajectories into
  stable reasoning patterns and transient noise
  (Theorem~\ref{thm:jdlg}).
\end{itemize}

All theoretical results are verified computationally on explicit
finite-dimensional examples, confirming that the spectral radius $\rho(B)$
accurately predicts convergence behavior.

\subsection*{Organization}
Section~\ref{sec:prelim} introduces notation, defines the trajectory space and
CoT operators, and recalls prerequisite results from semigroup theory.
Section~\ref{sec:main} contains our main theorems with complete proofs.
Section~\ref{sec:discussion} discusses implications, connections to related
work, and open questions. Section~\ref{sec:conclusion} summarizes our
contributions.
